\section{Continuations and exact residual spheres}\label{s:residual-spheres}

The block estimates test a continuation on many affine subspheres,
including subspheres of zero measure for the original uniform prior.
We first obtain Borel continuation functions on every such test. We
then determine the exact target law left by a Gaussian block.

\subsection{Borel rules and counterfactual continuations}

\begin{lemma}[Borel versions under the uniform prior]\label{s:borel}
Fix \(d\ge2\), a finite horizon, and finite state sets. For almost
every realization of the shared rule seed in Definition~\ref{s:model},
there are Borel transition and stopping kernels with the same state
sets and terminal output laws whose conditional joint law of the
uniform signal, state path, stopping index, and output is unchanged.
\end{lemma}
\begin{proof}
Fix a seed realization for which the conditional model requirements
hold. Such realizations have full seed probability.
For a fixed state and index, combine continuing and stopping choices
into a finite set of tagged destinations. Their probabilities form a
finite vector of completed-measurable nonnegative functions summing to
one on the pair space. Every completed-measurable real function has a
Borel version: approximate it by simple functions and replace their
countably many measurable level sets by Borel sets modulo null sets.
Choose such a version for each coordinate. On the Borel set where the
chosen coordinates are negative or fail to sum to one, replace the
vector by a fixed point mass. This set is null. Since there are only
finitely many indices and states, a single Borel
\(\gamma_d(dx)\,dy\)-null set \(Z\) contains every pair-input
modification.

Let \(S\) be uniform on \(S^{d-1}\) and pre-generate the independent
Gaussian rows. The joint independence in Definition~\ref{s:model}
ensures that, after fixing this seed, they retain this product law and
the initial state carries no information about them. For a fixed
\(x\ne0\), the coordinate of \(S\) in the
direction \(x/\norm x\) has a Lebesgue density when \(d\ge2\); this
will also follow from \eqref{s:coordinate-density} below. Scaling by
\(\norm x\) shows that the conditional law of \(\ip xS\) has a
density. The event \(x=0\) is Gaussian-null. Thus each pair
\((X_t,\ip{X_t}{S})\) has marginal law absolutely continuous with
respect to \(\gamma_d(dx)\,dy\) and avoids \(Z\) almost surely.
This conclusion remains valid along an adaptive state path, because
\[
 \Prob\{\text{current state}=v,\ (X_t,\ip{X_t}{S})\in Z\}
 \le \Prob\{(X_t,\ip{X_t}{S})\in Z\}=0.
\]
A finite union over the pre-generated indices is still null. Couple the
original and Borel finite kernels using the same fresh uniform random
numbers. Their choices agree outside \(Z\), so induction makes their
state paths and stopping indices identical almost surely. The terminal
output laws have no pair-input argument and can then be coupled
identically. This proves the joint-law assertion.
\end{proof}

The reduction preserves the averaged experiment; it need not preserve
the behavior at every exceptional individual signal. After making the
reduction, start the Borel learner at any specified state and index,
whether or not that state could have been reached. Its future success
probability \(G_v(s)\) is a Borel function of \(s\), by finite
composition and integration of the Borel kernels over fresh rows and
randomness. These are counterfactual continuations, rather than
conditional laws given a reached-state event.
Every bound below is uniform over the seed realizations covered by
the lemma, hence over almost every seed.
The original jointly measurable experiment supplies measurable
conditional success probabilities. We may therefore average those
original probabilities afterward; no jointly measurable choice of
Borel versions across seeds is needed.

For a block of \(k\) pre-generated rows, the complete transitions from
a fixed starting state give Borel destination probabilities
\(\kappa_v(X,y)\). If a padded program always has a destination, they
sum to one. The exact recursion is
\begin{equation}\label{s:recursion}
 G_{\rm in}(s)=
 \E_X\sum_v\kappa_v(X,Xs)G_v(s).
\end{equation}
All dependence of the block computation on its rows and labels remains
in \(\kappa_v\). The continuation uses only its state from the past.
Conditional on the shared seed, initialization is independent of the
signal and sample rows. It may be averaged at the end, after which the
same uniform estimates may be averaged over the seed in the original
jointly measurable experiment.

\subsection{Spherical coordinates and beta fractions}

We use uniform probability, never unnormalized surface area, on each
sphere. The following normalized formulas are classical; see Frankl
and Maehara \cite{FranklMaehara1990} for the geometric beta-law
context. We include their derivation and the conditional information
needed here.

\begin{lemma}[Coordinates of a uniform direction]\label{s:coordinates}
Let \(E=E_1\oplus E_2\) be an orthogonal decomposition with
\(\dim E=D\), \(\dim E_1=l\), and \(1\le l<D\). If
\(U\sim\sigma_E\), then
\[
 \norm{P_{E_1}U}^2\sim\Beta(l/2,(D-l)/2),
\]
and the two component directions are independent uniform directions,
independent of this fraction. In coordinates on \(E_1\), the
projection \(P_{E_1}U\) has density
\begin{equation}\label{s:coordinate-density}
 c(D,l)(1-\norm t^2)^{(D-l-2)/2}\ind_{\{\norm t<1\}},
 \qquad
 c(D,l)=\frac{\Gamma(D/2)}
              {\pi^{l/2}\Gamma((D-l)/2)}.
\end{equation}
Conditional on that projection, the complementary direction is
uniform on \(S(E_2)\). When \(l=0\) or \(l=D\), the squared
projection is respectively the point mass at \(0\) or at \(1\).
\end{lemma}
\begin{proof}
Write \(U=G/\norm G\) for a standard Gaussian \(G\) in \(E\).
The two Gaussian components are independent. Their squared lengths
divided by two are independent unit-scale gamma variables of shapes
\(l/2\) and \((D-l)/2\), and their directions are independent
uniform directions. For independent gamma variables \(A,B\) of
shapes \(\alpha,\beta>0\), the change of variables
\((A,B)=(qv,(1-q)v)\) has Jacobian \(v\). Their joint density therefore
factors into
\[
 \frac{q^{\alpha-1}(1-q)^{\beta-1}}{\mathrm B(\alpha,\beta)}
 \cdot
 \frac{v^{\alpha+\beta-1}e^{-v}}{\Gamma(\alpha+\beta)}.
\]
The fraction and sum are independent, and the fraction is
\(\Beta(\alpha,\beta)\). Polar coordinates on \(E_1\) give
\eqref{s:coordinate-density}; its normalizing constant follows from
the beta integral and the area
\(2\pi^{l/2}/\Gamma(l/2)\) of \(S^{l-1}\).
The independent component directions give the conditional assertion.
This proof includes \(D-l=1\), when \(S(E_2)\) has two points and the
coordinate-density exponent can be \(-1/2\). The endpoint
decompositions \(l=0,D\) are deterministic and require no beta law
with a zero shape.
\end{proof}

For \(D\ge3\), write \(c_D=c(D,1)\). The normalizing integrand
\((1-t^2)^{(D-3)/2}\) is at least \(1/2\) on
\(|t|\le D^{-1/2}\), so \(c_D\le\sqrt D\). For an exponent at least
one this lower bound follows from Bernoulli's inequality; the remaining
nonnegative exponents are immediate. This simple estimate will keep
all cap constants exponential in the ambient dimension.

\begin{lemma}[Beta moments and products]\label{s:beta-identities}
For \(A,B>0\) and \(Z\sim\Beta(A,B)\), a real \(q\ge0\) satisfies
\begin{equation}\label{s:beta-real}
 \E Z^{-q}
 =\frac{\mathrm B(A-q,B)}{\mathrm B(A,B)}
 =\frac{\Gamma(A-q)\Gamma(A+B)}
        {\Gamma(A)\Gamma(A+B-q)}
 \quad\text{if }q<A.
\end{equation}
The moment diverges when \(q\ge A\). If \(h\) is a nonnegative integer
with \(h<A\), then
\begin{equation}\label{s:beta-integer}
 \E Z^{-h}=\prod_{t=1}^h\frac{A+B-t}{A-t},
\end{equation}
where the product for \(h=0\) is one. Moreover, for \(C>0\), independent
\[
 Z\sim\Beta(A,B),\qquad X\sim\Beta(A+B,C)
\]
satisfy
\begin{equation}\label{s:beta-product}
 ZX\sim\Beta(A,B+C).
\end{equation}
\end{lemma}
\begin{proof}
Integrating \(z^{-q}\) against the beta density gives the first ratio
when the exponent \(A-q-1\) is greater than \(-1\). If it is at most
\(-1\), the integral diverges at zero. The identity
\(\mathrm B(x,y)=\Gamma(x)\Gamma(y)/\Gamma(x+y)\) gives
\eqref{s:beta-real}, and the gamma recurrence gives
\eqref{s:beta-integer}. For the product, take independent unit-scale
gamma variables \(G_1,G_2,G_3\) of shapes \(A,B,C\).
The fraction \(G_1/(G_1+G_2)\) is independent of \(G_1+G_2\), hence
of \((G_1+G_2)/(G_1+G_2+G_3)\). These two fractions have the stated
laws, and their product is \(G_1/(G_1+G_2+G_3)\), of law
\(\Beta(A,B+C)\).
\end{proof}

\subsection{What an exact block leaves}

Uniform orthogonal frames may be generated successively from uniform
sphere directions in perpendicular complements; Mezzadri gives an
authoritative exposition of this construction
\cite[Section~8, Theorem~4 and Corollary~1]{Mezzadri2007}.
Here is the short invariance fact we need. A square Gaussian matrix
\(A\) is invertible almost surely, and
\(Q=A(A^{\mathsf T}A)^{-1/2}\) is orthogonal. Left and right
orthogonal invariance of the Gaussian law show that \(Q\) is invariant
under both left and right multiplication by any fixed orthogonal
matrix. We call this the uniform orthogonal law. Its image of a fixed
frame is uniform. To see its conditional frame property, construct an
ordered orthonormal frame by choosing the first vector uniformly and
each successive vector uniformly in the perpendicular complement of
the previous ones. This joint law is left rotation invariant.
Uniqueness of invariant probability on any sphere, frame
space, or space of flags used below follows by averaging over \(Q\):
the average of a fixed object is independent of that object by
transitivity and right invariance, while averaging an invariant law
leaves it unchanged. Thus the successive construction has the same law
as the image under \(Q\), and its definition gives the claimed uniform
conditional law of the remaining frame.
All subspace averages below are measurable: if \(K=QK_0\) is represented
by an orthogonal matrix and a fixed subspace, then
\(\int f(Qu)\,d\sigma_{K_0}(u)\) is a Borel function of \(Q\) for
Borel \(f\), by parameter integration, and depends only on \(K\).

\begin{lemma}[Exact residual law]\label{s:residual}
Let \(s=z+r\omega\) with \(\omega\sim\sigma_H\), independently of
a matrix \(X\) with \(b\) independent standard Gaussian rows. Assume
\(z\perp H\), \(r>0\), and \(\norm z^2+r^2=1\).
Put \(D=\dim H\), and let \(1\le b<D\).
Let \(R_X\) be the row span of \(X\) restricted to \(H\), and set
\(K=H\cap R_X^\perp\). Conditional on \(X\), the joint experiment has
the representation
\begin{equation}\label{s:residual-representation}
 s=z+r\sqrt{1-t}\,e+r\sqrt t\,v,
 \quad
 t\sim\Beta((D-b)/2,b/2),\quad
 e\sim\sigma_{R_X},\quad v\sim\sigma_K,
\end{equation}
where the three draws are independent. The labels are
\[
 Xs=X\bigl(z+r\sqrt{1-t}\,e\bigr).
\]
Marginally, \(t\) is independent of \((e,K)\), \(e\) is uniform on
\(S(H)\), and, conditional on \(e\), \(K\) is uniform among the
\((D-b)\)-dimensional subspaces of \(H\cap e^\perp\).
\end{lemma}
\begin{proof}
The restriction of \(X\) to \(H\) has rank \(b\) almost surely and
its row span is rotation invariant. Conditional on \(X\),
Lemma~\ref{s:coordinates} applied to \(H=R_X\oplus K\) gives
\eqref{s:residual-representation}. In particular the residual
direction \(v\) is independent of \(t,e\) and remains uniform when
the full \(X,t,e\) are fixed. Since \(Xv=0\), the displayed label
identity follows. The restriction of \(X\) to \(R_X\) is injective,
so, given \(X\), the labels determine the projected component
\(r\sqrt{1-t}\,e\).

To reverse the order of the other draws, take a uniform orthonormal
basis of \(H\), let \(R_X\) be the span of its first \(b\) vectors,
and take its first vector for \(e\). By the invariance just proved,
this has the same law as a uniform row span and a uniform direction
in it. Conditional on \(e\), the remaining basis is a uniform frame
in \(H\cap e^\perp\). Its last \(D-b\) vectors span a uniform
\((D-b)\)-space there. The beta law of \(t\) does not depend on the
row span or its directions, which proves the claimed independence.
When \(b=1\), the conditional residual space is simply
\(H\cap e^\perp\); the statement still applies.
\end{proof}

\begin{figure}[H]
\centering
\begin{tikzpicture}[scale=1.25,>=Latex]
  \draw[thick] (0,0) circle (2);
  \draw[gray,dashed] (1.6,-1.2) -- (1.6,1.2);
  \draw[->,very thick,blue] (0,0) -- (1.6,0)
    node[midway,below=10pt] {\(p\in R_X\)};
  \draw[->,very thick,red] (1.6,0) -- (1.6,1.2)
    node[midway,right=5pt] {\(\rho v\in K\)};
  \draw[->,thick] (0,0) -- (1.6,1.2)
    node[midway,above left] {\(r\omega\)};
  \draw (1.42,0) -- (1.42,.18) -- (1.6,.18);
  \fill (0,0) circle (1.4pt) node[below left=3pt] {\(z\)};
  \fill (1.6,0) circle (1.4pt) node[below right=3pt] {\(z+p\)};
  \fill (1.6,1.2) circle (1.4pt) node[above right] {\(s\)};
  \fill (1.6,-1.2) circle (1.4pt);
  \node[align=center] at (-3.5,.7)
    {row-space projection\\fixed by the labels};
  \draw[->,gray] (-2.35,.5) to[bend right=15] (.65,-.04);
  \node[align=center] at (3.6,-.75)
    {uniform residual\\direction in \(K\)};
  \draw[->,gray] (2.8,-.6) to[bend left=15] (1.67,-.8);
\end{tikzpicture}

\caption{An orthogonal cross-section of the decomposition in
Lemma~\ref{s:residual}. The observed row-space component is
\(p=r\sqrt{1-t}\,e\); the remaining component has length
\(\rho=r\sqrt t\) and direction in \(K\).
The line meets the displayed circle in two points because this
cross-section is two-dimensional. In higher dimension the same
intersection is the residual subsphere \(z+p+\rho S(K)\).}
\label{s:residual-figure}
\end{figure}

Figure~\ref{s:residual-figure} shows the geometric content of the
conditioning. The subspaces in the lemma are analysis variables; the
learner is given exactly \(X,Xs\). The next consequence is the point
where the full data-dependent state choice is accounted for.

\begin{lemma}[Averaging after a block]\label{s:state-maximum}
In the setting of Lemma~\ref{s:residual}, let \(\mathcal V\) be nonempty and finite,
let \(g_v:S^{d-1}\to[0,1]\) be Borel, and let Borel weights
\(\kappa_v(X,y)\ge0\) satisfy \(\sum_v\kappa_v\le1\). Put
\[
 z'=z+r\sqrt{1-t}\,e,\qquad \rho=r\sqrt t,\qquad
 F_v(K)=\int g_v(z'+\rho u)\,d\sigma_K(u).
\]
Then
\begin{equation}\label{s:state-max-bound}
 \E_{X,s}\sum_v\kappa_v(X,Xs)g_v(s)
 \le \E_{t,e}\E_{K\mid e}\max_v F_v(K).
\end{equation}
The expectation on the left uses the independent uniform subsphere
target and Gaussian block.
\end{lemma}
\begin{proof}
Conditional on \(X,t,e\), the labels are \(Xz'\), the weights are
fixed, and \(u\) is uniform in \(K\). The conditional mean is
\(\sum_v\kappa_v(X,Xz')F_v(K)\), bounded by \(\max_vF_v(K)\).
This maximum no longer depends on the remaining matrix data except
through \(K\). The reversed-order part of Lemma~\ref{s:residual}
therefore gives the right side of \eqref{s:state-max-bound}.
\end{proof}

The maximum is taken before the matrix is replaced by its marginal
subspace law. It allows the selected state to depend arbitrarily on
the entire matrix and every exact label. No claim that a forward
posterior at a reached state is uniform on a subsphere is used.
